\documentclass{article}
\usepackage{calc}           % for \widthof
\usepackage{propositions}

%% Per-level dimension lists: every list dimension key accepts a
%% comma-separated list of values, one per nesting level.  The first applies
%% at level 1, the second at level 2, and so on; a gap falls back to the class
%% default for that level.  This file reports the dimensions the list actually
%% ends up with, so the expectations can be checked without measuring the page.

\newcommand\report{\texttt{LM=\the\leftmargin, LW=\the\labelwidth,
  LS=\the\labelsep}}

%% Print the width of some text (\widthof is only usable in a calc expression).
\newlength\testwd
\newcommand\wid[1]{\setlength\testwd{\widthof{#1}}\the\testwd}

\begin{document}

\section{Class defaults (baseline)}

\verb|\leftmargini| = \the\leftmargini, \verb|\leftmarginii| =
\the\leftmarginii, \verb|\leftmarginiii| = \the\leftmarginiii.

With no options set, each level should use the class default for its depth.

\begin{prop}
  \item Level 1: \report
  \begin{prop}
    \item Level 2: \report
    \begin{prop}
      \item Level 3: \report
    \end{prop}
  \end{prop}
\end{prop}

\section{A per-level list via \texttt{\textbackslash propoptions}}

\verb|leftmargin={3em, 1em}|: level 1 is 3em and level 2 is 1em.  Level 3 is
past the end of the list, which is treated like a gap: it falls back to the
class default (\the\leftmarginiii), rather than reusing the last value.

\begingroup
\propoptions{leftmargin={3em, 1em}}
\begin{prop}
  \item Level 1 (expect 3em): \report
  \begin{prop}
    \item Level 2 (expect 1em): \report
    \begin{prop}
      \item Level 3 (expect \the\leftmarginiii): \report
    \end{prop}
  \end{prop}
\end{prop}
\endgroup

\section{Gaps fall back to the class default}

\verb|leftmargin={, 0em}|: level 1 takes the class default (\the\leftmargini),
level 2 is 0em, so level-2 text aligns with level-1 text.

\begingroup
\propoptions{leftmargin={, 0em}}
\begin{prop}
  \item Level 1 (expect \the\leftmargini): \report
  \begin{prop}
    \item Level 2 (expect 0.0pt): \report
  \end{prop}
\end{prop}
\endgroup

\section{Per-level \texttt{labelwidth} sized to the widest label}

\verb|labelwidth={\widthof{(99)}, \widthof{a.}, \widthof{(iii)}}| gives each
level just enough room for its own widest label.

\begingroup
\propoptions{labelindent=0pt,
  labelwidth={\widthof{(99)}, \widthof{a.}, \widthof{(iii)}}}
\begin{prop}
  \item Level 1 (expect LW=\wid{(99)}): \report
  \begin{prop}
    \item Level 2 (expect LW=\wid{a.}): \report
    \begin{prop}
      \item Level 3 (expect LW=\wid{(iii)}): \report
    \end{prop}
  \end{prop}
\end{prop}
\endgroup

\section{Per-level \texttt{labelindent}}

\verb|labelindent=0em| at every level: the label starts at the enclosing
margin, so sub-item labels sit in the gap left by the outer label while the
sub-item text keeps the class-default indent.

\begingroup
\propoptions{labelindent=0em}
\begin{prop}
  \item Outermost item, labels flush with the enclosing margin.
  \begin{prop}
    \item Sub-item a: label in the gap, text indented.
    \item Sub-item b: same layout.
  \end{prop}
  \item Another outermost item, confirming dimensions are restored.
\end{prop}
\endgroup

The original mockup layout: \verb|labelindent={0em, -1.5em}| pushes level-2
labels left into the gap, and \verb|leftmargin={3em, 0em}| zeroes the level-2
margin so level-2 text aligns with level-1 text.

\begingroup
\propoptions{labelindent={0em, -1.5em}, leftmargin={3em, 0em}}
\begin{prop}
  \item Outermost item.  The level-two text below should align with this text.
  \begin{prop}
    \item Sub-item one: its text aligns with the level-one text above, and its
    label sits to the left, in the gap where the outer label was.
    \item Sub-item two: confirming alignment.
  \end{prop}
  \item Another outermost item, confirming level-one layout is unaffected.
\end{prop}
\endgroup

\section{A per-level list in an environment argument is env-local}

An argument styles only its own environment, so the nested list below falls
back to the class default rather than inheriting the second entry of the list.

\begin{prop}[leftmargin={3em, 1em}]
  \item Level 1 (expect 3em): \report
  \begin{prop}
    \item Level 2 (expect \the\leftmarginii, not 1em): \report
  \end{prop}
\end{prop}

With \verb|\propoptions| the same list does cascade, as in section~2 above.

\end{document}
